\section{Problem Formulation}
\label{sec:problem}

\subsection{Agent Trajectories and Addressable Evidence}

Let an agent execution be represented by an immutable trajectory
$\tau=(s_1,\ldots,s_n)$, where each span $s_i$ records an entity, operation,
parent relation, time interval, status, and a finite map of attributes. A
stable selector $p=(i,f)$ addresses field $f$ of span $s_i$. Resolving $p$
against $\tau$ yields an evidence item
\begin{equation}
    e=(i,f,v,H(v)),
\end{equation}
where $v$ is the canonical field value and $H$ is SHA-256 over its canonical
serialization. An evidence reference is valid only if the span and field
exist, the observed value equals $v$, and the recorded digest equals $H(v)$.
Validity establishes that a citation is authentic; it does not establish that
the evidence is relevant or sufficient for a causal claim.

\subsection{Structured Diagnosis}

A diagnosis report is
\begin{equation}
    d=(r,y,S,C,E,q,m),
\end{equation}
where $r\in\{\textsc{diagnosed},\textsc{no-failure},\textsc{abstained}\}$ is
the report status, $y$ is a failure type when $r=\textsc{diagnosed}$, $S$ is a set of
decisive spans, $C=(c_1,\ldots,c_k)$ is a bounded causal chain, $E$ is a set of
evidence references, $q\in[0,1]$ is a model-reported confidence, and $m$ is
provenance. Each claim $c_j$ names its causal stage and the subset of $E$ used
to support it. A report may instead return \textsc{no-failure} or abstain with
a structured reason. Contract~v1 does not encode a separate
responsible-entity field; attribution is represented by decisive spans and
causal claims. Entity-level evaluation is therefore defined only when the
reference label is derivable from span metadata or supplied through a separately
versioned contract. The contract constrains
syntax, identity, integrity, and state consistency; semantic support is
evaluated separately.

\subsection{Independent Verification Protocol}

Given $(\tau,d)$, verifier channel $k$ produces
\begin{equation}
    z_k=(u_k,F_k,G_k,m_k),
\end{equation}
where $u_k\in\{\textsc{verified},\textsc{needs-evidence},
\textsc{review-required}\}$, $F_k$ is a set of findings, $G_k$ is a set of
typed evidence gaps, and $m_k$ records verifier provenance. Hard deterministic
findings---for example, an invalid selector, hash mismatch, impossible time
order, or scope violation---make a report ineligible for automatic acceptance.
The semantic channel evaluates relevance, sufficiency, counter-evidence, and
alternative causal explanations. We call the protocol independently verified
when channel inputs and failure sources are deliberately separated and their
residual error correlation is measured; we do not assume statistical
independence a priori.

\subsection{Selective Diagnosis Risk}

Let $L^{\mathrm{rep}}(d,y^*)\in\{0,1\}$ be a report-level loss against an
execution-derived reference $y^*$. The loss is one when any prespecified
critical component is wrong: report status, failure type, decisive-span set,
causal relation, or evidence sufficiency. Entity attribution may be added only
for datasets whose gold responsibility label is derivable from the frozen span
representation. Equivalent sufficient-evidence sets and the adjudication
procedure are frozen before calibration. Adjudication is completed without
access to calibration scores or acceptance decisions, and every sampled report
receives a certification loss $\widetilde L^{\mathrm{rep}}$: it equals the
adjudicated loss when all critical labels resolve and is set conservatively to
one if any critical label remains unresolved. No report or group is excluded
because a critical label is unresolved.

The independent Bernoulli unit is a preregistered group $G_i$ keyed by the
tuple (task template, repository, environment); all report opportunities that
share the tuple belong to the same group. These groups are sampled independently
from the in-distribution group population. Their report opportunities
$j=1,\ldots,N_i$ and grouping keys are fixed without using calibration outcomes.
A frozen scalar error score $s(x)$ assigns larger values to riskier reports. At
candidate $\lambda$, report-level acceptance is
\begin{equation}
    B_{ij}(\lambda)=
    \mathbf{1}\{s(X_{ij})\leq\lambda\}\mathbf{1}\{H_{ij}=1\},
\end{equation}
where $H_{ij}$ is deterministic eligibility. The prespecified group policy
accepts a group iff it accepts at least one report and assigns a binary group
loss iff any accepted report has a frozen critical defect:
\begin{align}
    A_i^{G}(\lambda)
        &=\mathbf{1}\!\left\{\sum_{j=1}^{N_i}B_{ij}(\lambda)>0\right\},\\
    L_i^{G}(\lambda)
        &=\mathbf{1}\!\left\{\sum_{j=1}^{N_i}
          B_{ij}(\lambda)\widetilde L_{ij}^{\mathrm{rep}}>0\right\}.
\end{align}
The population group acceptance rate and primary selective risk are
\begin{align}
    a_G(\lambda) &= \Pr(A_i^G(\lambda)=1),\\
    p_\lambda^G &=
    \Pr(L_i^G(\lambda)=1\mid A_i^G(\lambda)=1).
\end{align}
Hard deterministic violations are never eligible, regardless of their learned
risk score.

For $n$ independently sampled calibration groups, define
\begin{equation}
    M_\lambda^G=\sum_{i=1}^{n}A_i^G(\lambda),\quad
    K_\lambda^G=\sum_{i=1}^{n}A_i^G(\lambda)L_i^G(\lambda).
\end{equation}
Empty group acceptance is infeasible, not zero risk, and
$K_\lambda^G/M_\lambda^G$ is defined only when $M_\lambda^G>0$. The analogous
individual-report ratio and its report counts are descriptive because reports
within a group need not be independent.

Before calibration, we freeze the score, deterministic tie-breaking, minimum
accepted-group count $m_{\min}$, and a finite candidate family $\Lambda$. For
overall failure probability $\delta$, let $U_\lambda$ be the one-sided
Clopper--Pearson upper bound for $p_\lambda^G$, using per-candidate failure
probability $\delta/|\Lambda|$ so that the bounds hold simultaneously over the
family~\cite{clopper1934use}. Explicitly,
\begin{equation}
    U_\lambda=
    \begin{cases}
        1, & M_\lambda^G=0\ \text{or}\ K_\lambda^G=M_\lambda^G,\\
        \begin{aligned}
        &\operatorname{Beta}^{-1}\!\left(
        1-\frac{\delta}{|\Lambda|}; K_\lambda^G+1,\right.\\[-1mm]
        &\qquad\left.M_\lambda^G-K_\lambda^G\right),
        \end{aligned} & \text{otherwise}.
    \end{cases}
\end{equation}
The feasible set is
\begin{equation}
    \mathcal{F}=\{\lambda\in\Lambda:
    U_\lambda\leq\alpha,\ M_\lambda^G\geq m_{\min}\}.
\end{equation}

\paragraph{Conditional finite-sample guarantee.}
For every fixed $\lambda\in\Lambda$, conditional on the number of accepted
groups, the erroneous accepted-group count is binomial with parameter
$p_\lambda^G$ under independent group sampling and a frozen pipeline. The
one-sided exact bounds therefore satisfy
\begin{equation}
    \Pr\!\left(\forall\lambda\in\Lambda:
    p_\lambda^G\leq U_\lambda\right)\geq 1-\delta
\end{equation}
by the union bound. Consequently, whenever $\mathcal{F}$ is nonempty, any
candidate selected from $\mathcal{F}$ by the frozen rule satisfies
$p_\lambda^G\leq\alpha$ with probability at least $1-\delta$. This statement
is conditional on the defined sampling population, frozen loss, complete
candidate family, and absence of calibration-dependent pipeline changes.

The chosen operating point maximizes calibration group acceptance over
$\mathcal{F}$, with the frozen tie-breaker, and is evaluated once on untouched
test data. If $\mathcal{F}$ is empty, the outcome is ``no feasible operating
point.'' This yields a calibration-only PAC certificate under frozen-pipeline
and i.i.d.-group assumptions~\cite{xu2025scrc}; it is not ordinary conformal
risk control applied directly to the random selective-risk ratio
~\cite{angelopoulos2024crc}. With zero erroneous accepted groups, one fixed
candidate, and one-sided 95\% confidence, the exact upper bound requires 59
accepted independent calibration groups for $\alpha=0.05$ and 29 for
$\alpha=0.10$.
Alternative treatments of unresolved labels are reported separately as
sensitivity analyses and do not inherit this certificate.

\paragraph{Untouched-test evaluation.}
The selected candidate is fixed before test access and evaluated exactly once.
Test counts and uncertainty estimate its realized performance; they neither
repeat candidate selection nor retroactively alter the calibration certificate.

\subsection{Bounded Evidence Acquisition}

When verification yields an actionable gap $G$, an acquisition policy
$\pi(G)$ may choose one action from a finite allowlist. Level~0 uses existing
evidence, Level~1 reveals an already recorded but omitted field or span, and
Level~2 replays at most one allowlisted read-only and idempotent operation. The
action, returned evidence, cost, and resulting diagnosis revision are logged.
Because acquisition changes the observed feature distribution, the decision
rule after acquisition is calibrated for the complete policy over routing,
action, diagnosis revision, re-verification, and final acceptance
~\cite{xu2026bcea}. If the acquisition policy is selected on calibration data,
the simultaneously corrected family is the frozen policy-by-threshold family;
otherwise the policy is frozen on development data before calibration. No
pre-acquisition threshold is reused after acquisition. If no action is
permitted or the post-acquisition report remains ineligible, the system
abstains.
